
% uncomment in non-corona semester
Hier haben wir euch die wichtigsten Termine der ersten Woche
zusammengefasst. Eine Einladung mit ausführlichen Beschreibungen der
Erstsemesterveranstaltungen solltet ihr schon mit eurem
Erstibrief\footnote{\url{https://teri.fsi.uni-tuebingen.de/anfibrief/}}
bekommen haben.\\
Die Anmeldung und mehr Informationen zu den Veranstaltungen findet man hier:

\link{https://eei.fsi.uni-tuebingen.de}{Anmeldeplattform für Fachschaftsveranstaltungen}

\definecolor{lightlightgray}{RGB}{235,235,235}
\newcommand{\event}{\cellcolor{lightlightgray}}

Der folgende Stundenplan gilt für den B. Sc. in Informatik. Euer eigener Stundenplan kann also von dem hier Gezeigten abweichen. Achtet also auf den Stundenplan in eurem Ersti-Brief. Die Ersti-Veranstaltungen sowie die Mathe und Informatik-Vorlesungen sollten aber alle betreffen. 

%\textbf{Woche 1 - (08. - 14. April)}\\
%\\
%\begin{tabular}{|l|p{0.13\linewidth}|p{0.13\linewidth}|p{0.13\linewidth}|p{0.13\linewidth}|p{0.13\linewidth}|p{0.13\linewidth}|} \hline
% Zeit & Montag & Dienstag & Mittwoch & Donnerstag & Freitag & Samstag \\ 
% \hline \hline
% 08 -- 09 & & & & & &  \\ \cline{1-1} 
% 09 -- 10 & & & & & &  \\ \cline{1-1}
% 10 -- 11 & & & & & &  \\ \cline{1-1}
% 11 -- 12 & & & & & &  \\ \cline{1-1}
% 12 -- 13 & & & & & &  \\ \cline{1-1} 
% 13 -- 14 & & & & & &  \\ \cline{1-1}
% 14 -- 15 & & & & & &  \\ \cline{1-1}
% 15 -- 16 & & & & & &  \\ \cline{1-1} \cline{6-6}
% 16 -- 17 & & & & &\scriptsize Begrüßung &  \\ \cline{1-1} 
% 17 -- 18 & & & & & \scriptsize Fachbereich &  \\ \cline{1-1} \cline{6-6}	
% 18 -- 19 & & & & & &   \\ \cline{1-1}
% 19 -- 20 & & & &\cellcolor{lightlightgray} \footnotesize{Spieleabend} & \cellcolor{lightlightgray}&   \\ \cline{1-1}
% 20 -- 21 & & & &\cellcolor{lightlightgray} & \cellcolor{lightlightgray}&  \\ \cline{1-1}
% 21 -- 22 & & & &\cellcolor{lightlightgray} & \cellcolor{lightlightgray}& \\ \hline
% \end{tabular}
%\\
%{\scriptsize TBA = to be announced. Wird auf Fachschaftswebsite\footnote{\url{https://eei.fsi.uni-tuebingen.de}}  bekanntgegeben.} \\
%{\scriptsize Grau unterlegt = Veranstaltungen der Fachschaften Informatik und Kognitionswissenschaft }
%\vfill

% ------------------ Wintersemester ------------------------
%\pagebreak
%\textbf{Erste Vorlesungswoche (13. - 18. Oktober)}\\
%\\
%\begin{tabular}{|l|p{0.13\linewidth}|p{0.13\linewidth}|p{0.13\linewidth}|p{0.13\linewidth}|p{0.13\linewidth}|p{0.13\linewidth}|} \hline
%Zeit & Montag & Dienstag & Mittwoch & Donnerstag & Freitag & Samstag \\ 
%\hline \hline \cline{3-3}
%08 -- 09 & \footnotesize{Mathe 1} & & \footnotesize{Mathe 1} & & &\\ \cline{1-1}
%09 -- 10 &  & & & & &  \\ \cline{1-1} \cline{2-2} \cline{4-4}
%10 -- 11 & \footnotesize{Tech. Info 1} & & & & & \\ \cline{1-1}
%11 -- 12 & & & & & & \\ \cline{1-1} \cline{2-2}
%12 -- 13 & & & & & & \\ \cline{1-1} \cline{4-4}
%13 -- 14 & & & \footnotesize{Tech. Info 1}  & & & \\ \cline{1-1} \cline{3-3} \cline{4-4} \cline{5-5}
%14 -- 15 & &\footnotesize{Prakt. Info 1} & &\footnotesize{Prakt. Info 1} & & \\ \cline{1-1}
%15 -- 16 & & & & & & \\ \cline {1-1} \cline{3-3} \cline{5-5}
%16 -- 17 & & & & & & \\ \cline{1-1}
%17 -- 18 & & & & & & \\ \cline{1-1}
%18 -- 19 & & & & & &  \\ \cline{1-1}
%19 -- 20 & &\cellcolor{lightlightgray} \footnotesize{Filmeabend} &\cellcolor{lightlightgray} \footnotesize{Spieleabend} & &\cellcolor{lightlightgray} \footnotesize{Kneipentour} & \cellcolor{lightlightgray} \footnotesize{Grillen} \\ \cline{1-1}
%20 -- 21 & &\cellcolor{lightlightgray} &\cellcolor{lightlightgray}  & & \cellcolor{lightlightgray} &\cellcolor{lightlightgray}\\ \cline{1-1}
%21 -- 22 & &\cellcolor{lightlightgray} & \cellcolor{lightlightgray} & &  \cellcolor{lightlightgray} &\cellcolor{lightlightgray}\\ \hline
%\end{tabular}
%\\
% {\scriptsize TBA = to be announced. Wird auf Fachschaftswebsite\footnote{\url{https://eei.fsi.uni-tuebingen.de}}  bekanntgegeben.}
% {\scriptsize Grau unterlegt = Veranstaltungen der Fachschaften Informatik und Kognitionswissenschaft }

% ----------  Sommersemester  --------------
% \setlength{\arrayrulewidth}{0.3mm}  % Setzt die Dicke der Linien

%\pagebreak
 \textbf{Erste Vorlesungswoche (13. - 18. April)}\\
 \\
 \begin{tabular}{|l|p{0.13\linewidth}|p{0.13\linewidth}|p{0.13\linewidth}|p{0.13\linewidth}|p{0.13\linewidth}|p{0.13\linewidth}|} \hline
  Zeit & Montag & Dienstag & Mittwoch & Donnerstag & Freitag & Samstag \\ 
  \hline \hline 
  08 -- 09 & & & & & &\\ \cline{1-1}
  09 -- 10 &  & & & & &  \\ \cline{1-2} \cline{4-4} 
  10 -- 11 & \footnotesize{Mathe 2} & &  \footnotesize{Mathe 2} & & & \\ \cline{1-1}
  11 -- 12 & & & & & & \\ \cline{1-1} \cline{2-2} \cline{4-4}
  12 -- 13 & & & & & & \\ \cline{1-1} 
  13 -- 14 & & & & & & \\ \cline{1-1} \cline{3-3} \cline{5-5}
  14 -- 15 & &\footnotesize{Prakt. Info 2} & &\footnotesize{Prakt. Info 2} & & \\ \cline{1-1}
  15 -- 16 & & & & & & \\ \cline {1-3} \cline{5-5}
  16 -- 17 & \footnotesize{Tech. Info 2} & & & & & \\ \cline{1-1}
  17 -- 18 & & & & & & \\ \cline{1-1}
  18 -- 19 & & & & & &  \\ \cline{1-3} \cline{6-6}
  19 -- 20 & & \cellcolor{lightlightgray} \footnotesize{Spieleabend} & & & \cellcolor{lightlightgray} \footnotesize{Kneipentour} & \\ \cline{1-1}
  20 -- 21 & & \cellcolor{lightlightgray} & & & \cellcolor{lightlightgray} & \\ \cline{1-1}
  21 -- 22 & & \cellcolor{lightlightgray} & & & \cellcolor{lightlightgray} & \\ \hline
 \end{tabular}
 \\
%{\scriptsize TBA = to be announced. Wird auf Fachschaftswebsite\footnote{\url{https://eei.fsi.uni-tuebingen.de}}  bekanntgegeben.}
%{\scriptsize Grau unterlegt = Veranstaltungen der Fachschaften Informatik und Kognitionswissenschaft }

\normalsize
Außerdem gibt es noch viele andere Ersti-Veranstaltungen, die ihr auf unserer Website oder im Ersti-Brief nachlesen könnt. 
%Am \textbf{Sonntag 20. Oktober} gibt es zusätzlich noch Capture the Flag auf dem Sand. Außerdem gibt es eine weitere Wanderung am \textbf{Samstag 26. Oktober} und das Clubhausfest am \textbf{Donnerstag 31. Oktober}.

%Die diesjährige Ersti-Hütte findet von \textbf{Freitag, 01.November bis Sonntag, 03.November} statt. 

\pagebreak

% ----------  Sommersemester  --------------
\textbf{Mathe 2}\\[4mm]
Die Vorlesung\footnote{vollständiger Name: Mathematik für Informatik 2: Lineare Algebra} wird dieses Semester von Prof. Britta Dorn gehalten und umfasst in der Regel die drei Themenblöcke Algebraische Strukturen (Gruppen, Ringe, Körper), Vektorräume \& lineare Abbildungen (unter anderem viel Rechnen mit Vektoren und Matrizen) und Euklidische \& unitäre Räume (Skalarprodukte und noch mehr Rechnen mit Vektoren und Matrizen).\\
Die Vorlesung findet dieses Semester auf der Morgenstelle in den Hörsälen N6 (Montag) und N2 (Mittwoch) statt. Zusätzlich zur Vorlesung gibt es auch ein Tutorium.\\

\textbf{Praktische Info 2}\\[4mm]
Diese Vorlesung\footnote{vollständiger Name: Praktische Informatik 2: Imperative und objektorientierte Programmierung} wird dieses Semester von Dr. Philipp Schuster gehalten und ist eine gründliche Einführung in die Entwicklung von Programmen und die dazugehörigen formalen Grundlagen.\\
Man verbringt einen Großteil der Zeit damit, eigene Programme zu schreiben und wird
dabei Schritt für Schritt an Konzepte wie Rekursion oder Abstraktion herangeführt.\\
Die Vorlesung findet dieses Semester auf der Morgenstelle im Hörsaal N7 statt. Zusätzlich zur Vorlesung gibt es hier ebenfalls ein Tutorium.\\

\textbf{Technische Info 2}\\[4mm]
Diese Vorlesung\footnote{vollständiger Name: Technische Informatik 2: Informatik der Systeme} (von den älteren Studierenden oft IdS genannt) wird von Prof. Michael Menth gehalten und gibt einen Überblick über verschiedenste Themen der Technischen Informatik. Behandelt werden unter anderem Internet, Kodierung, Assemblerprogrammierung, Rechnerarchitektur, Betriebssysteme und Energieversorgung.\\
Die Vorlesung findet dieses Semester auf der Morgenstelle im Hörsaal N3 statt. Statt wie in einer normalen Vorlesung wird der Vorlesungsstoff bei dieser Vorlesung oft über Videos vermittelt, in der Präsenzstunde (VDP-Stunde) werden dann offene Fragen geklärt, aktuelle Themen nochmal aufgegriffen und manchmal auch potenzielle Klausuraufgaben besprochen.\\

\textbf{übK}\\[4mm]
Im Studienplan wird empfohlen, 3 ECTS an übK\footnote{überfachlich berufsfeldorientierte Kompetenzen, manchmal auch Liberal Education genannt} zu belegen. Hier lässt sich fast alles belegen, was mit ECTS-Punkten versehen ist, benotet wird und kein Sport ist.

\newpage
\normalsize

